\section{Introduction}
\label{sec:motivation}

\emph{Automatic item generation} (AIG) can reduce the workload of producing new test items and their answer keys. In template-based AIG, an author writes a single template, or \emph{item model}, from which a computer generates many items, each carrying its own answer key \citep{bejar2002generative,gierl2013aig}. AIG has been applied mainly to items built from text. Items whose response is a \emph{picture} have been harder to reach. In K--12 mathematics, the picture that matters most is the geometry diagram.

Producing a correct diagram is harder than it appears. Consider a simple request:

\begin{quote}
\emph{``Draw an acute triangle $ABC$ with $\angle A = 60^\circ$ and $\angle B = 70^\circ$, then draw the altitude from $C$, meeting $AB$ at $H$.''}
\end{quote}

To be correct, the drawing must satisfy several requirements \emph{exactly}. Angle $A$ is $60^\circ$. The point $H$ lies on the segment $AB$, not past its end. The altitude meets $AB$ at a true right angle. A model that draws directly has no step at which the drawing is checked against these requirements. Whether they hold is a matter of chance. When they fail, they fail silently.

The same exactness makes such diagrams hard to score. Most geometry benchmarks pose the reverse task. They show the model a finished diagram and ask a question about it \citep{lu2021intergps,lu2024mathvista}. They do not ask whether a model can \emph{produce} a correct figure. The benchmarks that do score generated pictures were built for everyday images. They rely on a model-based judge that inspects the output \citep{huang2025t2icompbenchpp,wang2025genexam}. A judge can report that a shape \emph{resembles} a right triangle. It cannot certify that an angle is exactly $90^\circ$. For geometry diagrams, then, neither producing the item nor scoring it is straightforward.

Our approach removes both the guessing and the inspection. The model describes the construction in words: ``$H$ is the foot of the perpendicular from $C$ to $AB$.'' A geometry engine then computes the exact coordinates. The model decides \emph{what} to build. The engine determines \emph{where} every point lies. This single change addresses both problems.

\begin{itemize}
\item \textbf{Better drawings.} A solver performs the computation, so a right angle is exactly right rather than approximately so.
\item \textbf{Automatic scoring.} The finished figure is a set of exact coordinates, so each requirement can be checked directly. The question ``is this angle $90^\circ$?'' replaces a judge's impression of the picture.
\end{itemize}

That is this paper's contribution: a \emph{machine-verifiable item-model framework} for geometry diagrams. One item model emits both the prompt and a \emph{computed} answer key. The key is a machine-checkable list of geometric requirements. The figure produced in response is verified against that key. We make two empirical claims and are deliberate about their strength. Structured construction improves how often models get the geometry right (\cref{sec:benchmark}). Human validation shows the automatic score runs optimistic, so it is not yet a validated high-stakes score (\cref{sec:scoring}). A fuller leaderboard and cost analysis appear in a companion paper. Here the benchmark is evidence for these two claims, not the contribution itself.

The rest of the paper describes how the model describes a diagram (\cref{sec:pipeline}), how scoring follows from that same description (\cref{sec:scoring}), what we found on six large models (\cref{sec:benchmark}), and why this matters for educational measurement (\cref{sec:measurement}).
